% ATLAS report design v1. Synthetic mockup; not a generated strategy report.
% Compiled with LuaHBTeX 1.24.0 / TeX Live 2026 on Windows, 13 September 2026.
\documentclass[11pt,a4paper]{article}
\usepackage[margin=20mm,headheight=16pt]{geometry}
\usepackage{fontspec}
\setmainfont{TeX Gyre Heros}[Numbers=Monospaced]
\setmonofont{Latin Modern Mono}[Scale=0.9]
\usepackage[spanish]{babel}
\usepackage[table]{xcolor}
\usepackage{booktabs,tabularx,longtable,fancyhdr,hyperref,pgfplots}
\pgfplotsset{compat=1.18}
\definecolor{ink}{HTML}{272B2D}
\definecolor{copper}{HTML}{975435}
\definecolor{cream}{HTML}{F4F1EA}
\definecolor{ivory}{HTML}{FFFCF6}
\definecolor{slate}{HTML}{526B80}
\definecolor{rulegray}{HTML}{DCD7CF}
\color{ink}
\hypersetup{colorlinks=true,urlcolor=copper,linkcolor=copper,pdftitle={ATLAS: maqueta sintética}}
\setlength{\parindent}{0pt}
\setlength{\parskip}{7pt}
\renewcommand{\arraystretch}{1.3}
\linespread{1.15}
\pagestyle{fancy}
\fancyhf{}
\fancyhead[L]{\small\textbf{ATLAS}\quad Investigación cuantitativa}
\fancyhead[R]{\small Maqueta v1 · Datos sintéticos}
\fancyfoot[L]{\small atlas-report-template-v1}
\fancyfoot[R]{\small\thepage}
\renewcommand{\headrulewidth}{0.3pt}
\newcommand{\AtlasSection}[1]{\par\addvspace{12pt}{\large\bfseries\color{copper}#1}\par\nobreak\vspace{3pt}}
\begin{document}
{\small\color{copper}\textbf{INFORME DE DESARROLLO}}\par
{\LARGE\bfseries Tendencia y preservación de capital}\par
ACTIVO DEMO · EUR\hfill 02/01/2024 — 08/01/2024\par
\colorbox{cream}{\parbox{\dimexpr\linewidth-2\fboxsep}{\small
\textbf{Maqueta con datos sintéticos.} Cinco observaciones para revisar el diseño.
No representa un backtest ejecutado ni una estrategia validada.}}
\AtlasSection{Resumen del periodo}
\begin{tabularx}{\linewidth}{@{}Xrrr@{}}
\toprule
Referencia & Patrimonio EUR & Variación & Caída máxima\\\midrule
Ejemplo sintético & 10.125,00 & 1,25\,\% & 0,30\,\%\\
Referencia sintética & 10.080,00 & 0,80\,\% & 0,20\,\%\\
Efectivo & 10.000,00 & 0,00\,\% & 0,00\,\%\\\bottomrule
\end{tabularx}
\AtlasSection{Evolución del patrimonio}
\begin{tikzpicture}
\begin{axis}[width=.94\linewidth,height=60mm,axis lines=left,
  axis line style={rulegray},tick style={rulegray},grid=major,
  grid style={rulegray!45},ylabel={EUR},xtick={1,2,3,4,5},
  xticklabels={02/01,03/01,04/01,05/01,08/01},scaled y ticks=false,
  yticklabel={\pgfmathprintnumber[fixed,precision=0,1000 sep={.},dec sep={,},assume math mode=true]{\tick}},
  tick label style={font=\small\sffamily},
  legend style={draw=none,font=\small,at={(0.5,1.05)},anchor=south,legend columns=2}]
\addplot[color=slate,thick] coordinates {(1,10000)(2,10100)(3,10069.7)(4,10100)(5,10125)};
\addlegendentry{Ejemplo sintético}
\addplot[color=ink,dashed] coordinates {(1,10000)(2,10050)(3,10029.9)(4,10050)(5,10080)};
\addlegendentry{Referencia sintética}
\end{axis}
\end{tikzpicture}
\AtlasSection{Lectura y límites}
La muestra ilustra alineación de cifras, jerarquía y contraste. No debe inferirse
significancia estadística a partir de cinco puntos. El gráfico es vectorial y sus
series se distinguen por trazo además de color. La reserva queda fuera del informe.
\newpage
\AtlasSection{Contexto de cálculo}
\begin{tabularx}{\linewidth}{@{}lX@{}}
\toprule
Campo & Contenido de la maqueta\\\midrule
Origen & Datos sintéticos diseñados para composición\\
Rango solicitado & 02/01/2024 — 08/01/2024, extremos inclusivos\\
Rango efectivo & Cinco sesiones ficticias\\
Política & Demostración visual; sin motor económico\\
Estado de evidencia & Sintético, no acreditado\\\bottomrule
\end{tabularx}
\AtlasSection{Detalle y trazabilidad}
Los informes reales incluirán aquí sus operaciones o posiciones, parámetros,
comisiones, supuestos, versiones y huellas, sin trasladar las métricas de toda
la historia a un subperiodo. Las tablas extensas repetirán cabeceras.
\begin{longtable}{@{}p{0.2\linewidth}p{0.72\linewidth}@{}}
\toprule Identificador & Valor\\\midrule\endfirsthead
\toprule Identificador & Valor (continuación)\\\midrule\endhead
Plantilla & \texttt{atlas-report-template-v1}\\
Tipo & Maqueta sintética\\
Fuente & Este documento; no hay proveedor de mercado\\\bottomrule
\end{longtable}
\AtlasSection{Reproducción}
El paquete conserva fuente TeX, JSON, CSV, recursos locales y manifiesto.
Esta maqueta se ha compilado con LuaHBTeX 1.24.0 de TeX Live 2026 en Windows,
con sockets y shell-escape desactivados. Sus cifras siguen siendo sintéticas.
\end{document}
